\documentclass[11pt,a4paper]{article}
\usepackage{amsmath,amssymb,amsthm}
\usepackage{graphicx}
\usepackage[margin=2.2cm]{geometry}
\usepackage{xcolor}
\usepackage{titlesec}
\usepackage{fancyhdr}
\usepackage{booktabs}
\usepackage{enumitem}
\usepackage{float}
\usepackage{hyperref}
\hypersetup{colorlinks=true,linkcolor=blue!70!black,urlcolor=blue!70!black}

\renewcommand{\topfraction}{0.9}
\renewcommand{\bottomfraction}{0.9}
\renewcommand{\textfraction}{0.05}
\renewcommand{\floatpagefraction}{0.5}
\setcounter{topnumber}{4}
\setcounter{bottomnumber}{4}
\setcounter{totalnumber}{6}
\raggedbottom

\definecolor{primary}{RGB}{0,51,102}
\definecolor{accent}{RGB}{204,102,0}

\titleformat{\section}{\Large\bfseries\color{primary}}{\thesection}{1em}{}
\titleformat{\subsection}{\large\bfseries\color{primary!80!black}}{\thesubsection}{1em}{}

\pagestyle{fancy}
\fancyhf{}
\fancyhead[L]{\small\color{primary}Essential Summary}
\fancyhead[R]{\small\color{primary}Li et al.\ (2013)}
\fancyfoot[C]{\thepage}
\renewcommand{\headrulewidth}{0.5pt}

\begin{document}

\begin{center}
{\LARGE\bfseries\color{primary}Essential Summary}\\[12pt]
{\large\bfseries Generalized Pad\'{e} Approximation Method\\for Homoclinic Orbits of Strongly Nonlinear\\Autonomous Oscillators}\\[8pt]
{\normalsize Zhenbo Li, Jiashi Tang, Ping Cai}\\[4pt]
{\small\itshape Chinese Journal of Theoretical and Applied Mechanics\\Vol.\ 45, No.\ 3, 461--464 (2013)\quad DOI: 10.6052/0459-1879-12-277}\\[6pt]
{\small\rule{0.7\textwidth}{0.5pt}}\\[4pt]
{\small\color{accent}This is an author-prepared essential summary, not the original paper.\\Please cite the published version.}
\end{center}

\vspace{0.5em}

\section{Background and Motivation}

The determination of homoclinic orbits is central to understanding global bifurcations, chaos, and solitary wave solutions in nonlinear dynamics \cite{Nayfeh1995}. Existing Pad\'{e}-based approaches share a common limitation: the numerator and denominator of the approximant are restricted to polynomial functions, which limits flexibility in capturing the exponential decay behavior characteristic of homoclinic orbits.

This paper introduces the \textbf{Generalized Pad\'{e} Approximation (GPA) method}, a theoretical extension that removes this restriction by allowing the approximant to be constructed from \emph{any} kind of function series. This provides a unified framework that unifies all existing Pad\'{e}-type modifications as special cases.

\section{Core Innovation: The Generalized Pad\'{e} Approximation}

\subsection{From Classical to Generalized}

The classical Pad\'{e} approximant of a formal power series $f(z)=\sum_{i=0}^{\infty} a_i z^i$ is a rational function satisfying:
\begin{equation}
f(z) - \frac{P_L(z)}{Q_M(z)} = \mathcal{O}(z^{L+M+1}), \qquad Q_M(0)=1.
\end{equation}
The generalized definition extends $H_m$ to $\widehat{H}_m$:
\begin{equation}
\boxed{\widehat{H}_m = \Big\{P : P(z)=\sum_{i=0}^m a_i g_i(z)^i,\; g_i:\mathbb{C}\to\mathbb{C},\; a_i\in\mathbb{C}\Big\}},
\end{equation}
and $G(m,n) = \{R = P/Q : P\in\widehat{H}_m,\; Q\in\widehat{H}_n\setminus\{0\}\}$. The \textbf{GPA[$L/M$]$_f$} is then the rational function $P_L/Q_M\in G(L,M)$ satisfying the same $\mathcal{O}(z^{L+M+1})$ condition. When $g_i(z)=z^i$, the GPA reduces to the classical Pad\'{e} approximant.

\subsection{Homoclinic Orbit Construction via Hyperbolic Functions}

For $\ddot{x} + g(x) = 0$, the power series solution $x(t) = \sum_{i=0}^{\infty} a_i t^i$ is first obtained. For a homoclinic orbit through saddle $H(h,0)$, the initial amplitude $a$ and initial velocity ($a_1=0$) are determined by the energy integral $\int_h^a g(x)\,dx = 0$.

Since $\lim_{t\to\pm\infty} x(t) = h$, a power series cannot capture the asymptotic behavior. The key innovation --- construct a GPA using \textbf{hyperbolic functions} that naturally satisfy the boundary conditions:
\begin{equation}
\boxed{\text{GPA}[L/M] = \frac{\displaystyle\sum_{i=0}^{L} \alpha_i \operatorname{sech}^i t}{\displaystyle 1 + \sum_{i=1}^{M} \beta_i \operatorname{sech}^i t}}.
\end{equation}

The $L+M+1$ unknown coefficients $\{\alpha_i, \beta_i\}$ are determined by expanding the GPA into a Taylor series around $t=0$ and matching coefficients with the power series $\{a_i\}$. The boundary condition $x(\pm\infty)=h$ is automatically satisfied ($\operatorname{sech}^i t \to 0$). The coefficient matching reduces to a system of $L+M$ linear equations coupled with one nonlinear condition from the energy integral.

\subsection{Advantages}

The $\operatorname{sech}^i t$ basis offers three advantages: (1) its Maclaurin series is simpler than exponentials, reducing algebraic effort when matching terms; (2) it satisfies $x(\pm\infty)=h$ automatically; (3) it applies to \emph{any} $g(x)$ that is a polynomial or rational function.

\section{Key Results: Three Numerical Examples}

\subsection{Example 1: Strongly Asymmetric Helmholtz--Duffing}
\begin{equation}
\ddot{x} + c_1 x + c_2 x^2 = 0,\qquad c_1=-3,\; c_2=100.
\end{equation}
Despite the enormous asymmetric nonlinearity ($c_2=100\gg 1$), GPA[$2/2$] matches Runge--Kutta with excellent agreement:
\begin{equation}
\text{GPA}[2/2] = \frac{0.005\operatorname{sech} t + 0.07\operatorname{sech}^2 t}{1 + 0.611111\operatorname{sech} t + 0.055555\operatorname{sech}^2 t}.
\end{equation}

\begin{figure}[htbp]
\centering
\includegraphics[width=0.75\textwidth]{p7-fig1-homoclinic.png}
\caption{\textbf{Figure 1.} (Original Fig.\ 1) Homoclinic orbit of the Helmholtz--Duffing oscillator with $c_1=-3, c_2=100$. Solid: Runge--Kutta. Dotted: GPA[$2/2$]. The method maintains accuracy despite the strongly asymmetric nonlinearity.}
\end{figure}

\subsection{Example 2: Full Cubic--Quadratic Helmholtz--Duffing}
\begin{equation}
\ddot{x} + c_1 x + c_2 x^2 + c_3 x^3 = 0,\qquad c_1=-2,\; c_2=5,\; c_3=15.
\end{equation}
Both left and right homoclinic orbits are obtained simultaneously with $L=M=2$, demonstrating the method handles asymmetric cubic--quadratic nonlinearities without modification.

\begin{figure}[htbp]
\centering
\includegraphics[width=0.75\textwidth]{p7-fig2-homoclinic.png}
\caption{\textbf{Figure 2.} (Original Fig.\ 2) Homoclinic orbits for $c_1=-2, c_2=5, c_3=15$. Solid: Runge--Kutta. Dotted: GPA[$2/2$], capturing both left and right homoclinic orbits.}
\end{figure}

\subsection{Example 3: Duffing--Harmonic Oscillator (Rational Restoring Force)}
\begin{equation}
\ddot{x} + \frac{\lambda x + \mu x^3}{1 + \nu x^2} = 0,\qquad \lambda=-1,\; \mu=1,\; \nu=1.
\end{equation}
With $L=M=3$, the GPA successfully captures homoclinic orbits of this rational-restoring-force oscillator. The exact solution cannot be expressed in elementary or classical elliptic functions --- yet the GPA handles it without modification.

\begin{figure}[htbp]
\centering
\includegraphics[width=0.75\textwidth]{p7-fig3-duffing-harmonic.png}
\caption{\textbf{Figure 3.} (Original Fig.\ 3) Homoclinic orbits of the Duffing-harmonic oscillator with $\lambda=-1, \mu=1, \nu=1$. Solid: Runge--Kutta. Dotted: GPA[$3/3$]. The method handles rational restoring forces without additional computational complexity.}
\end{figure}

\section{Significance and Impact}

\begin{enumerate}[leftmargin=2em]
\item \textbf{Theoretical unification:} The GPA definition unifies all existing Pad\'{e}-type modifications (quasi-Pad\'{e}, exponential-function approximants, modified Pad\'{e}) under a single theoretical framework.

\item \textbf{Computational efficiency:} The $\operatorname{sech}^i t$ basis produces simpler Taylor expansions than exponential-based alternatives, reducing the cost of coefficient matching at the same approximation order.

\item \textbf{Broad class coverage:} The method handles polynomial $g(x)$ of arbitrary degree and rational $g(x)$ --- a significant advantage over methods restricted to specific restoring force forms.
\end{enumerate}

\section{Limitations}

\begin{itemize}[leftmargin=2em]
\item Conservative systems only ($\varepsilon=0$); damped oscillators require separate methods.
\item Approximation order $L=M$ chosen empirically; no a priori error estimate.
\item Single-degree-of-freedom systems only.
\end{itemize}

\section{Connection to the GHFP/MGHFP Series}

This paper (together with its English version, Chin.\ Phys.\ B 2014) and the GHFP/MGHFP series represent two \textbf{complementary methodological families}. Both aim at analytical treatment of strongly nonlinear oscillators, but via different mathematical routes: GPA constructs homoclinic/heteroclinic solutions directly through rational functions of hyperbolic functions in the time domain; GHFP/MGHFP uses perturbation theory in the angular domain to derive amplitude--parameter relationships. The $\operatorname{sech}^i t$ and $\tanh^i t$ basis functions introduced here also appeared in the GHFP framework (JSV 2013), highlighting the shared mathematical intuition across both families.

\vfill
\begin{center}
\small\rule{0.5\textwidth}{0.3pt}\\[4pt]
\small\color{primary} Author: Zhenbo Li (lizhenbo@usc.edu.cn)\\
\small University of South China, School of Mathematics and Physics\\
\small Hunan Key Laboratory of Mathematical Modeling and Scientific Computing
\end{center}

\begin{thebibliography}{2}
\bibitem{Nayfeh1995}
A.H.\ Nayfeh and B.\ Balachandran, \textit{Applied Nonlinear Dynamics}, Wiley, 1995.
\bibitem{P8}
Z.\ Li, J.\ Tang, and P.\ Cai, ``A generalized Pad\'e approximation method of solving homoclinic and heteroclinic orbits of strongly nonlinear autonomous oscillators,'' \textit{Chin.\ Phys.\ B}, vol.\ 23, no.\ 12, p.\ 120501, 2014.
\end{thebibliography}

\end{document}
